\section{Introduction}
\label{sec:intro}

Speech-driven facial animation aims to generate mouth motion temporally synchronized with speech, supporting digital avatars, game production, and embodied interaction. Recent work has steadily improved the visual quality of talking faces across portrait animation~\citep{ji2021audio,zhang2023sadtalker,xu2024vasa,jiang2024loopy,xu2024hallo}, 3D facial animation~\citep{fan2022faceformer,xing2023codetalker}, and audio-conditioned neural rendering~\citep{guo2021ad,chen2025echomimic,ma2026esgaussianface}. Compact facial representations such as 3D morphable models, FLAME, and blendshape parameterizations provide practical control spaces~\citep{blanz2023morphable,li2017learning,menzel2022automated}, among which ARKit defines semantically meaningful facial action coefficients~\citep{apple_arkit} and has become a standard interface between speech and mouth motion in production. However, how a model decides which mouth shape a given speech sound requires is still left to be discovered from acoustics.

The difficulty lies in the size of the mapping that must be learned from the waveform alone. For instance, a 10-second utterance at 30\,Hz spans 300 frames over 32 mouth-related ARKit channels, so a model must emit 9{,}600 coefficients whose correct values are governed by only about one hundred phonemes. Existing audio-driven methods regress these coefficients from learned acoustic representations~\citep{cudeiro2019voca,richard2021meshtalk,fan2022faceformer,park2023said}, which asks the model to rediscover phoneme identity, phoneme boundaries, and the articulatory action of each phoneme as a by-product of fitting motion. Such implicit learning weakens local correspondence in long utterances, and leaves the model without any notion that a bilabial closure and a wide jaw opening are different articulatory events instead of two nearby points in coefficient space.

Existing attempts to inject linguistic structure into this mapping fall into two lines. Early efforts build viseme inventories and phoneme-to-shape rules~\citep{cohen1993modeling,edwards2016jali,taylor2017deep,zhou2018visemenet}, which make articulation explicit but compress speech into a small discrete inventory, losing the coarticulation that determines how strongly each action is realized. More recently, a line of work injects textual or linguistic cues as auxiliary features of an audio-driven regressor~\citep{fan2022joint,xu2024kmtalk}, and language-aware facial models~\citep{kang2026semanticface,wu2025keyframeface} show that language-aligned supervision helps estimate facial actions. Although these methods enrich the input, the cue is consumed as one more feature vector and the model is never told what action a phoneme category names. In essence, both lines leave the connection between a phonetic unit and a named facial control outside the learning objective. Fortunately, recent multimodal language models~\citep{xu2025qwen3omni} consume audio and text jointly and follow natural-language instructions, which makes it possible to state that connection instead of learning it.

In this paper, we aim to advance speech-driven mouth animation from implicit acoustic regression to explicit articulation-grounded generation, where the model receives the linguistic content, the phonetic units, and the articulatory meaning of those units as inputs, and produces facial controls carrying the same names. Two obstacles stand in the way, namely \emph{boundary mismatch} and \emph{value imbalance}. On the one hand, grounding is only meaningful when a phoneme and the motion it produces occupy the same time span, which long unsegmented recordings do not provide. On the other hand, coefficient values concentrate near zero because most mouth channels stay inactive at any instant, so a token-level objective spends most of its capacity predicting inactivity and under-supervises the large values that shape visible motion. To this end, we propose AudioMouth, which formulates mouth animation as schema-constrained sequence generation with a multimodal language model. Specifically, to resolve boundary mismatch, AudioMouth constructs semantically aligned audio-linguistic units in which one interval selects the audio, transcript, phonemes, and target motion at once, as illustrated in Figure~\ref{fig:overview}. To resolve value imbalance, it applies \emph{distribution-balanced supervision} that downweights frequent near-zero targets so active channels receive proportionally stronger signal, and it generates the coefficient sequence under explicit \emph{articulation grounding} that names the control each phoneme group drives. Since next-token likelihood measures neither how a trajectory moves nor what shape it produces, AudioMouth finally refines the generator with Group Relative Policy Optimization, using a reward over coefficient accuracy, temporal dynamics, and a \emph{simplified lip-geometry discrepancy} (SLMD) that reads 3D mouth deformation out of a fixed blendshape basis without rendering or landmark detection.

Experiments on a subject-disjoint split of BEAT~\citep{liu2022beat} show that AudioMouth improves every reported metric over four Audio2ARKit baselines, reducing Fr\'echet distance by 57\% and coefficient MSE by 70\% relative to the strongest baseline. In summary, our key contributions are three-fold.
\begin{itemize}[leftmargin=1em, itemsep=-1ex, topsep=0ex]
    \item \textbf{Articulation-Grounded Generator:} To the best of our knowledge, AudioMouth is the first framework that generates ARKit mouth trajectories inside a multimodal language model, where phonemes and their articulatory meaning are explicit inputs and the named facial controls are the output schema.
    \item \textbf{Motion-Aware Refinement:} We propose a sequence-level refinement stage whose reward evaluates coefficient accuracy, temporal dynamics, and mouth geometry, and we introduce SLMD to obtain 3D geometric feedback directly from coefficient differences at negligible cost.
    \item \textbf{Strong Performance:} AudioMouth outperforms four Audio2ARKit baselines on all five metrics under a subject-disjoint split, and our analyses attribute the gain to articulation grounding for distributional fidelity, balanced supervision for motion amplitude, and geometric feedback for across-script consistency.
\end{itemize}
